\chapter{System Design}

\section{Overview}

This chapter describes the concrete design of the AI powered quiz and
coding challenge platform. It builds directly on
the technology choices justified in the Technology Review, demonstrating
how those choices were applied to produce a coherent, working system. The
chapter is organised from coarse to fine grain: the overall architecture
is presented first, followed by detailed descriptions of each major
subsystem, the database schema, the machine learning training pipelines,
the frontend architecture, and the deployment model. Diagrams are used
throughout to augment the written explanation.

The system consists of three containerised services orchestrated by Docker
Compose \cite{merkel2014docker,dockercompose2025}: a PostgreSQL 16
database, a Python FastAPI backend, and a React and TypeScript frontend
served by Nginx \cite{nginx2025}. These three services are loosely coupled
through well defined interfaces. The backend exposes a REST API
\cite{fielding2000rest} conforming to the OpenAPI 3.0 specification
\cite{openapi2025}, and the frontend communicates exclusively through
that API using JSON \cite{json2014} over HTTP. No service depends on the
internal state of any other service, and the only inter service contracts
are the REST API and the SQL schema.

\begin{figure}[H]
    \centering
    \includegraphics[width=\linewidth]{./images/design/high_level_system_architecture.png}
    \caption{High level system architecture showing the three Docker
    Compose services and their communication paths.}
    \label{fig:system-architecture}
\end{figure}

\section{Docker Compose Infrastructure}

The deployment unit of the project is a Docker Compose file declaring
three services: \texttt{db}, \texttt{backend}, and \texttt{frontend}.
Each runs in an isolated container but communicates with the others
over Docker's internal bridge network. Environment variables such as
database credentials, the OpenAI API key, and the JWT secret are
injected at runtime from a \texttt{.env} file and never committed to
version control.

The \texttt{db} service runs the official \texttt{postgres:16-alpine}
image with a named Docker volume \texttt{pgdata} mounted at
\texttt{/var/lib/postgresql/data} for persistence across container
restarts. Health checks via \texttt{pg\_isready} combined with
\texttt{depends\_on: condition: service\_healthy} ensure the backend
does not start before the database is ready. The backend container is
built from a custom Dockerfile based on \texttt{python:3.11-slim} with
Node.js 22 installed alongside to support the JavaScript and TypeScript
runtimes required by the code execution engine. The two serialised
model artefacts, \texttt{topic\_model.pkl} and
\texttt{difficulty\_model.pkl}, are copied in at build time. The
application is served by Uvicorn \cite{lathkar2023} listening on port
8000. The frontend Dockerfile uses a two stage build: in the first,
\texttt{node:22-alpine} runs \texttt{npm run build} to produce a Vite
\cite{vite2025} static bundle; in the second, the compiled assets are
copied into an \texttt{nginx:stable-alpine} image whose
\texttt{nginx.conf} handles single page application fallback routing,
gzip compression, and one year cache headers on content hashed assets.

\section{Backend Architecture}

The backend is a FastAPI application \cite{fastapi2025} whose entry
point is \texttt{main.py}. FastAPI's \texttt{lifespan} context manager
is used to run \texttt{init\_db()} at startup, which creates the four
database tables if they do not already exist, using idempotent
\texttt{CREATE TABLE IF NOT EXISTS} and
\texttt{ALTER TABLE ... ADD COLUMN IF NOT EXISTS} statements so that
migrations on an existing database are safe. The application is
structured around four concerns: AI model integration, machine learning
inference, code execution, and user management, which includes
authentication, gamification, and social features.

\begin{figure}[H]
    \centering
    \includegraphics[width=\linewidth]{./images/design/backend_dependency_architecture.png}
    \caption{Backend module dependency diagram showing call relationships
    between the FastAPI entry point, service modules, AI parsers, and
    machine learning inference wrappers.}
    \label{fig:backend-deps}
\end{figure}

\subsection{API Routes}

The FastAPI application exposes the following groups of HTTP endpoints.
All request and response bodies are validated against Pydantic
\texttt{BaseModel} subclasses defined in \texttt{pydantic\_models.py}
\cite{pydantic2025}, which means that malformed or type incorrect input
is rejected at the API boundary with a descriptive 422 error before it
reaches any business logic.

\begin{itemize}
    \item \texttt{POST /prompt}: accepts a \texttt{PromptRequest}, which
    bundles the user prompt, the quiz type, the target language, and the
    number of questions; invokes the AI generation pipeline; and returns
    a list of enriched \texttt{QuizSchema} or
    \texttt{CodingQuestionSchema} objects.
    \item \texttt{POST /run-code}: accepts source code and a language
    identifier, executes the code in a sandboxed subprocess, and returns
    stdout and stderr.
    \item \texttt{POST /submit-code}: accepts source code, a language
    identifier, and a JSON array of test cases; runs the code against
    each test case; and returns a per case pass or fail result.
    \item \texttt{POST /model}: toggles the active AI backend between
    OpenAI and Ollama.
    \item \texttt{POST /register} and \texttt{POST /login}: user
    registration and authentication, returning a signed JWT on success.
    \item \texttt{GET /verify-email} and \texttt{POST /resend-verification}:
    support the email verification flow that gates first time login.
    \item \texttt{GET /me}, \texttt{PATCH /me/profile}: retrieve and
    partially update the authenticated user's profile fields.
    \item \texttt{POST /add-xp}, \texttt{POST /save-quiz-result}: award
    experience points and persist a completed quiz result.
    \item \texttt{GET /user-stats}, \texttt{GET /quiz-history}: retrieve
    aggregated statistics and recent quiz sessions for the authenticated
    user.
    \item \texttt{GET /users/search}, \texttt{GET /users/\{id\}/profile},
    \texttt{GET /users/\{id\}/stats}: search for users and retrieve public
    profiles and friends only statistics.
    \item \texttt{POST /friends/request/\{id\}},
    \texttt{POST /friends/respond},
    \texttt{GET /friends/requests},
    \texttt{GET /friends/requests/count},
    \texttt{GET /friends},
    \texttt{GET /friends/count},
    \texttt{DELETE /friends/\{id\}}: the full social friendship lifecycle.
    \item \texttt{GET /}, \texttt{GET /health}: root and server status
    utility endpoints.
\end{itemize}

\subsection{AI Generation Pipeline}

The AI generation pipeline is the most complex subsystem in the backend.
The active AI backend is tracked by a global \texttt{current\_model}
variable in \texttt{ai\_models.py}, toggled via \texttt{POST /model}.
Both backends use \texttt{httpx.AsyncClient} with a 120 second timeout
and await the full response before parsing begins, which simplifies the
parser logic at the cost of not being able to show incremental output.

Once the raw response arrives, the appropriate parser
(\texttt{parser\_openai.py} or \texttt{parser\_ollama.py}) extracts the
JSON content from the model's reply. For MCQ quizzes, each question
object is passed to \texttt{predict\_tags\_for\_question(title,
question)} to obtain predicted topic tags. For coding quizzes,
\path{compute\_coding\_metadata()} is called first to extract numeric
and lexical features (character count, word count, constraint language,
code tokens), and these features are passed to
\texttt{predict\_difficulty\_for\_question()} alongside the question's
AI generated metadata to obtain a difficulty label. Both enrichments
are attached to the schema object before it joins the output list. The
Ollama parser additionally strips option prefix characters such as
\texttt{"A: "} from MCQ options, and the OpenAI backend wraps parsing
in a three attempt retry on \texttt{json.JSONDecodeError},
\texttt{KeyError}, or \texttt{ValueError}.

\subsection{Machine Learning Inference Pipeline}

\textbf{Model loading.} Both machine learning models are loaded into
memory as in process singletons. They are deserialised from their
\texttt{.pkl} artefacts once and held in module level variables for the
lifetime of the Uvicorn process. The difficulty model
(\texttt{difficulty\_model.pkl}) is loaded eagerly at module import
time in \texttt{difficulty\_service.py}. The \texttt{\_load\_model()}
function is called at module level, outside any function body, so the
\texttt{joblib.load} \cite{joblib2025} call executes when Uvicorn first
imports the module during startup. The topic model
(\texttt{topic\_model.pkl}) is loaded lazily: \texttt{\_load\_model()}
is called inside \texttt{predict\_topics()} and returns immediately if
the module level \texttt{\_pipeline} variable is already populated.
This means the topic model is deserialised on the first quiz generation
request after startup, then retained in memory for all subsequent
requests. Both strategies ensure that the inference cost per request is
limited to a single \texttt{pipeline.predict\_proba()} call and a
threshold comparison, which is negligible compared to the AI generation
latency.

A compatibility fix is applied when loading the difficulty model: the
helper functions defined in the \texttt{\_\_main\_\_} namespace at
training time are re injected into \texttt{sys.modules['\_\_main\_\_']}
before \texttt{joblib.load} is called, preventing an
\path{AttributeError} during deserialisation.

\textbf{Topic tag prediction.} The topic tag classifier is a scikit
learn \texttt{Pipeline} \cite{pedregosa2011,buitinck2013api} wrapping a
\texttt{FeatureUnion} of two \texttt{TfidfVectorizer} instances (word
level unigram to trigram and character level three to five n grams)
followed by a \texttt{OneVsRestClassifier} wrapping a
\texttt{LogisticRegression} estimator with balanced class weights. At
inference time, the \path{predict\_topic\_from\_parts(title,
description, top\_n=5)} function concatenates the title and description
with a period separator, calls \texttt{pipeline.predict\_proba()} to
obtain a probability score for each of the trained labels, applies the
per label optimised thresholds (stored in the \texttt{.pkl} bundle
alongside the pipeline and the \texttt{MultiLabelBinarizer}), and
returns up to five tags sorted by descending probability. If no label
exceeds its threshold, the label with the highest probability is
returned anyway, guaranteeing that every generated question carries at
least one meaningful topic annotation.

\textbf{Difficulty prediction.} The difficulty classifier is a scikit
learn \texttt{Pipeline} wrapping a \texttt{ColumnTransformer} that
applies separate preprocessing paths to numeric and text features,
followed by the best performing ensemble estimator selected during
training. The numeric path applies \texttt{SimpleImputer(strategy='median')}
followed by \texttt{StandardScaler} to the numeric features derived
from submission statistics and problem description analysis. The text
path applies a \texttt{FeatureUnion} of word level and character level
\texttt{TfidfVectorizer} instances to the raw description text. At
inference time, \texttt{predict\_difficulty\_for\_question(schema)}
receives a \texttt{CodingQuestionSchema} object, maps its fields to the
column names the model expects, calls
\texttt{extract\_description\_features()} to derive the text and keyword
features from the question body, assembles a single row pandas DataFrame
\cite{mckinney2010pandas}, and calls \texttt{pipeline.predict()} to
obtain one of the three labels: easy, medium, or hard. If the model is
unavailable, for example if the \texttt{.pkl} file is missing from the
image, the function falls back to returning \texttt{medium}.

\subsection{Difficulty Classifier Training Pipeline}

The difficulty classifier training pipeline, implemented in
\path{ml_models/difficulty_classifier/difficulty_predictor.py},
merits a more detailed description. The dataset is constructed by
iterating over the CodeNet \cite{puri2021codenet} per problem submission
CSV files and aggregating per problem statistics (mean CPU time, mean
memory, mean accepted code size, and acceptance rate). Problems with
fewer than ten submissions are dropped because their aggregates are
statistically meaningless. Problem descriptions are extracted from the
corresponding HTML files and parsed to produce a plain text form, a
word count, a character count, a count of numeric constraint tokens,
and a set of binary keyword flags indicating the presence of
algorithmic vocabulary such as recursion, dynamic programming, graph,
sort, search, string, math, greedy, stack, and complexity.

Labels are derived from acceptance rate by percentile binning: the top
tercile by acceptance rate is labelled easy, the middle tercile medium,
and the bottom tercile hard. This produces a balanced three class
dataset by construction. The resulting training matrix contains thirty
four numeric features (base statistics, log transforms, ratios, and
interaction terms) plus one text feature that is subsequently expanded
into word and character n grams by the preprocessor. The split uses
stratified sampling to produce a sixty, twenty, twenty train, validation,
and test partition. Seven candidate models are trained using
\texttt{RandomizedSearchCV} \cite{bergstra2012random} with forty
iterations and five fold cross validation, scoring on F1 macro. The
best model, XGBoost \cite{chen2016xgboost} in the final run, is
retrained on the combined train and validation sets and evaluated once
on the held out test set. The full pipeline, including the
\texttt{LabelEncoder}, is serialised to \texttt{difficulty\_model.pkl} via joblib.

\subsection{Topic Tag Classifier Training Pipeline}

The topic tag training pipeline, in \texttt{ml\_models/tag\_classifier/train.py},
follows a similar structure but is adapted to the multi label nature of
the problem. The LeetCode \cite{leetcode2025} CSV is loaded, rows with
null descriptions or null tag lists are dropped, and the tag strings
are parsed from their stored representation into lists. Tags appearing
in fewer than sixty problems are filtered out, and any row with no
remaining tags is dropped. A \texttt{MultiLabelBinarizer} converts the
tag lists to binary indicator vectors and is saved alongside the model
for use at inference.

The training matrix is a combined \path{cleaned\_title + '. ' +
cleaned\_description} string per problem. The feature pipeline is a
\texttt{FeatureUnion} of a word level TF IDF vectoriser (up to fifty
thousand features, unigrams through trigrams) and a character level
TF IDF vectoriser (up to twenty thousand features, three to five
character n grams). Four candidate architectures are trained using
\texttt{GridSearchCV} with three fold cross validation, scoring on F1
micro: logistic regression with word features only, logistic regression
with combined word and character features, an SGD classifier with word
features, and a classifier chain \cite{read2011classifierchains} of
logistic regressions with combined features.

A two stage threshold optimisation is then performed on the validation
set. Stage A finds a single global optimum threshold by sweeping from
0.15 to 0.70 in steps of 0.01 and choosing the value that maximises F1
micro. Stage B computes a per label optimum by sweeping each label
independently from 0.10 to 0.80 in steps of 0.025. The final per label
threshold is a blend of the per label optimum and the global optimum
with a shrinkage factor of 0.5. This shrinkage is a conservative bet
against overfitting: pure per label optimisation on a validation set
with small support per label is vulnerable to noise, whereas pure
global optimisation discards the benefit of knowing that rare labels
need lower thresholds. The 0.5 blend was chosen because it exactly
halves the influence of the per label sample on the final threshold,
which is a principled starting point when there is no stronger prior.
The best model is then retrained on train plus validation combined and
evaluated once on the test set, and the pipeline, the binarizer, the
thresholds, and the class list are all serialised together in
\texttt{topic\_model.pkl}.

\subsection{Code Execution Engine}

The code execution engine is implemented in
\texttt{services/code\_executor.py} and runs inside the FastAPI process
rather than as a separate service. For both \texttt{POST /run-code} and
\texttt{POST /submit-code}, the engine writes the user's source code to
a temporary file using Python's \texttt{tempfile} module, spawns the
appropriate runtime as a child process using \texttt{subprocess.run()},
captures stdout and stderr, and deletes the temporary file in a
\texttt{finally} block before returning. A five second timeout is
enforced on every subprocess call. Submissions that exceed this limit
are terminated and a timeout message is returned to the caller.

The runtime invocation strategy varies by language. Python uses the
\texttt{python} interpreter directly. JavaScript uses \texttt{node}.
TypeScript is executed directly using Node.js 22 with the
\texttt{--experimental-strip-types} flag, which removes type annotations
at runtime and eliminates the need for a separate transpilation step.
Java requires a two step process: \texttt{javac} compiles the source file
to a \texttt{.class} file, then \texttt{java} executes it. C\# uses the
\texttt{dotnet} CLI. Node.js 22 is installed in the backend Docker image
to support both JavaScript and TypeScript execution. The Python
interpreter is provided by the base image.

For test case evaluation (\texttt{POST /submit-code}), the engine
iterates over the list of test cases provided in the request body. For
Python submissions, the function name is extracted via a regex, then a
test harness is appended that calls the function with the supplied
input arguments and prints a marked result line to stdout. The harness
output is parsed and compared to the expected value. The full list of
per case results is returned to the frontend for rendering.

\subsection{Authentication and User Management}

Authentication follows the stateless JWT pattern standardised in IETF
RFC 7519 \cite{jwt2015}. When a user registers, their password is
hashed using bcrypt \cite{provos1999bcrypt} with a generated salt
before storage. On login, the stored hash is compared to the provided
password using \texttt{bcrypt.checkpw()}. If it matches, a JWT is
generated using PyJWT with the user's id as the subject claim and a
twenty four hour expiry encoded in the \texttt{exp} claim. The token
is signed using HMAC SHA256 with a secret loaded from the environment.
The frontend stores the token in \texttt{localStorage} and includes it
in the \texttt{Authorization: Bearer <token>} header on all
authenticated requests. Backend route handlers that require
authentication call \texttt{verify\_token()} from \texttt{core/auth.py},
which decodes and validates the JWT and returns the user ID for use
in the handler.

Email verification is enforced at registration time. A secure random
token generated by \texttt{secrets.token\_urlsafe(32)} is stored on the
user row together with a timestamp. A verification email is sent via
SMTP (with TLS), and the user cannot log in until the link in that
email is clicked. The verification endpoint only accepts tokens younger
than twenty four hours, a rule enforced at the SQL level using
\texttt{verification\_sent\_at > NOW() - INTERVAL '24 hours'}.

The experience point and levelling system is implemented in
\texttt{services/users.py}. When a quiz result is submitted via
\texttt{POST /save-quiz-result}, experience is awarded based on the
number of correct answers. The threshold for levelling up follows the
formula \texttt{xp\_for\_level = 100 * level}, so a level 1 user needs
100 experience points to reach level 2, a level 2 user needs 200 to
reach level 3, and so on. When accumulated experience meets or exceeds
the threshold, the level is incremented and the experience pool is
reduced by the threshold amount, allowing carry over. A \texttt{while}
loop handles the rare case where a single grant is large enough to
cover multiple levels at once.

\section{Database Design}

\begin{figure}[H]
    \centering
    \includegraphics[width=\linewidth]{./images/design/database_entity_diagram.png}
    \caption{Entity relationship diagram for the four database tables. The users table is the central entity, with one to many relationships to both quiz\_results and quiz\_sessions. The friendships table models a directed self referential relationship on users, with a unique constraint preventing duplicate requests and a check constraint preventing self friendship.}
    \label{fig:database-design}
\end{figure}

The database schema consists of four tables: \texttt{users},
\texttt{quiz\_results}, \texttt{friendships}, and \texttt{quiz\_sessions}.
The schema is created at application startup by \texttt{init\_db()}
using raw SQL \texttt{CREATE TABLE IF NOT EXISTS} statements, meaning
the tables are created automatically on first run without requiring a
separate migration tool.

\subsection{Users Table}

The \texttt{users} table is the central entity of the schema. In
addition to authentication fields (\texttt{email}, \texttt{password\_hash}),
it stores profile data (\texttt{display\_name}, \texttt{bio},
\texttt{avatar\_url}), gamification state (\texttt{exp}, \texttt{level}),
and email verification state (\texttt{email\_verified},
\texttt{verification\_token}, \texttt{verification\_sent\_at},
\texttt{verified\_at}). The \texttt{email} column carries a
\texttt{UNIQUE} constraint, enforcing that no two accounts share an
email address at the database level, as a complement to the application
level duplicate check.

\subsection{Quiz Results Table}

The \texttt{quiz\_results} table records the outcome of completed quiz
sessions. Notably, the full AI generated quiz content is not persisted
in this table. It exists only in React component state for the duration
of the session. Only the metadata necessary for the user's history and
statistics is written to the database: the prompt that generated the
quiz, the quiz type, the score, the predicted topic tags (stored as a
PostgreSQL \texttt{TEXT[]} native array column), the programming
language if applicable, and the completion timestamp. This design
deliberately keeps the database lean and avoids the cost of serialising
and storing potentially large JSON documents per session. Where a full
record of the session is needed, the \texttt{quiz\_sessions} table
stores the complete question set as a \texttt{JSONB} column, and this
supports the recent history view.

\subsection{Friendships Table}

The \texttt{friendships} table models the social graph as a directed
edge list. Each row represents a friendship request from
\texttt{requester\_id} to \texttt{addressee\_id} with a status of
either \texttt{pending} or \texttt{accepted}. Two database level
constraints prevent invalid states: a \texttt{UNIQUE(requester\_id,
addressee\_id)} constraint prevents duplicate requests between the same
pair of users, and a \texttt{CHECK (requester\_id != addressee\_id)}
constraint prevents self friendship. A confirmed mutual friendship is
represented by a single row with \texttt{status = 'accepted'}. The
application layer treats both users as friends regardless of which
direction the original request was sent.

\section{Frontend Architecture}

The frontend is a single page application built with React
\cite{react2013} and TypeScript \cite{typescript2025}, compiled by Vite
\cite{vite2025}. At runtime the browser loads a single HTML document
(\texttt{index.html}) containing script tags pointing to the Vite
generated JavaScript bundles. React Router handles client side
navigation, mapping URL paths to page components without triggering
full page reloads. Nginx's \texttt{try\_files \$uri /index.html}
directive ensures that direct URL access or page refresh always returns
the root HTML document so that React Router can take control.

\subsection{Route Structure}

The application exposes eleven routes, defined in \texttt{App.tsx}. Six
of these routes are protected by an authentication guard that redirects
unauthenticated users to the login page: \texttt{/prompt},
\texttt{/quiz}, \texttt{/code-sandbox}, \texttt{/profile},
\texttt{/social}, and \texttt{/stats}. The remaining five routes are
publicly accessible: \texttt{/} (landing), \texttt{/login},
\texttt{/register}, \texttt{/about}, and \texttt{/user/:userId} (public
profile). This separation reflects the application's requirement that
content generation and social features are gated behind an account,
while public profiles and informational pages are accessible without
authentication.

\subsection{State Management}

Global application state is managed by a single hand rolled
\texttt{AuthContext} built from React's \texttt{createContext} and
\texttt{useState} primitives. The context exposes five values:
\texttt{user} (the authenticated user object), \texttt{token} (the raw
JWT string), \path{isAuthenticated}, \texttt{isLoading}, and two
mutation functions \texttt{login} and \texttt{logout}. On login, both
the user object and the token are written to \texttt{localStorage} so
that the session survives a page refresh. On page load,
\texttt{AuthContext} reads from \texttt{localStorage} to rehydrate the
session without requiring a new login for the current browser tab.

All other state, quiz content, editor content, social list data, is
local component state managed with \texttt{useState} hooks inside the
relevant page component. This approach avoids the complexity of a
global state management library for a codebase of this size, and the
clear separation between global authentication state and local UI
state makes the data flow straightforward to reason about.

\subsection{Key Page Components}

\texttt{PromptPage} is the primary entry point to quiz generation. It
renders the \texttt{PromptForm} component, which collects the user's
topic prompt, quiz type (MCQ or coding), programming language (for
coding quizzes), and number of questions. On submission, it posts to
\texttt{POST /prompt} and awaits the complete enriched quiz payload.
On success, it navigates to either \texttt{/quiz} or
\texttt{/code-sandbox}, passing the quiz data via React Router's
location state.

\texttt{QuizPage} receives the MCQ quiz data from location state and
renders the questions one at a time, tracking the user's selected
answer and score in local state. Session state is also persisted
to \texttt{sessionStorage} under a key derived from the user id and
session id, so that a refresh of the page does not lose in progress
answers.

\texttt{CodeSandboxPage} renders an interactive coding environment for
the generated coding challenges, built on top of the Monaco Editor
\cite{monaco2025}. The editor, language selector, test case panel, and
results panel are all managed as local state within this component.
Submissions are sent to \texttt{POST /submit-code} and the per case
pass or fail results are rendered inline next to the expected versus
actual output.

\texttt{SocialPage} handles the full friends lifecycle: searching for
users by display name, sending friend requests, accepting or rejecting
incoming requests, and viewing the current friends list.

\section{End to End Data Flow}

\begin{figure}[H]
    \centering
    \includegraphics[width=\linewidth]{./images/design/end_to_end_data_flow.png}
    \caption{End to end data flow for a quiz generation session, showing the request path from browser through JWT authentication, AI generation via OpenAI or Ollama, and ML enrichment. MCQ questions receive a predicted difficulty label only, while coding challenges receive both difficulty and topic tags. Both models are loaded at startup as in process singletons. The session concludes with results persisted and experience points awarded.}
    \label{fig:end-to-end}
\end{figure}

To make the interaction between components concrete, this section
traces the complete path of a quiz generation and submission session.
The authenticated user submits a form at \texttt{/prompt} requesting
five Python coding challenges on a chosen topic. The browser sends
\texttt{POST /prompt} with the bearer JWT in the request header, and
FastAPI validates the body against the \texttt{PromptRequest} Pydantic
model before any handler logic runs. \texttt{ai\_models.py} prepends
the coding prompt template to the user's topic and dispatches the full
prompt to the active backend via \texttt{httpx.AsyncClient} with a 120
second timeout. When the response arrives, the parser extracts the
array of question objects and enriches each one: topic tags are
obtained from \texttt{predict\_tags\_for\_question()} and a difficulty
label from \texttt{predict\_difficulty\_for\_question()}, both of
which consult their in memory model singleton. The enriched
\texttt{CodingQuestionSchema} list is serialised and returned, and the
frontend stores it in local state before navigating to
\texttt{/code-sandbox}. Submissions from the sandbox are sent to
\texttt{POST /submit-code}, executed against the stored test cases by
the code executor, and the pass or fail result list is rendered inline.
On session completion, \texttt{POST /save-quiz-result} writes a
\texttt{quiz\_results} row and awards experience points, levelling up
the user if the threshold is crossed.

\section{Deployment Architecture}

The system is deployed locally using Docker Compose on a single host.
Containers communicate over Docker's default bridge network by service
name (the backend resolves the database at hostname \texttt{db}, and
Nginx proxies requests to \texttt{backend:8000}). The containerised
architecture is inherently portable: the same \texttt{docker-compose.yml}
could run on any Linux host provisioned on AWS EC2, Google Compute
Engine, or DigitalOcean with no code changes. A production deployment
would additionally require TLS termination (Nginx with Let's Encrypt),
a process supervisor, and a managed database service to decouple
database lifecycle from application containers.

\section{Design Decisions and Standards Compliance}

The machine learning inference pipeline follows the scikit learn
\texttt{Pipeline} convention \cite{buitinck2013api}, which ensures that
all preprocessing steps are encapsulated in a single serialisable
object. This prevents training and serving skew: the same
\texttt{TfidfVectorizer} vocabulary and IDF weights, and the same
\texttt{StandardScaler} parameters that were fitted on the training
data, are applied at inference time without any manual re
implementation. Separating the training code from the inference code
(the training scripts live in \texttt{ml\_models/}, the inference
wrappers live in \texttt{backend/ml/}) is a standard pattern for
productionising machine learning models and reflects the principle
that the artefact produced by training is a self contained, deployable
unit. The REST API conforms to the OpenAPI 3.0 specification
\cite{openapi2025}, authentication uses JWT (RFC 7519 \cite{jwt2015})
with HMAC SHA256 signing, and password storage uses bcrypt
\cite{provos1999bcrypt}, consistent with OWASP guidance
\cite{owasp2021}.